\begin{table}[htbp]
\centering
\caption{FL Price Premium by Modality and Minimum-Bidder Constraint}
\label{tab:modal_id}
\begin{threeparttable}
\small
\begin{tabular}{lcccc}
\toprule
 & (1) & (2) & (3) & (4) \\
 & Preg\~ao & Convite & Convite & Convite \\
 & (no rule) & (all) & ($n_{\text{gen}} \geq 3$) & ($n_{\text{gen}} < 3$) \\
 &  &  & voluntary & binding \\
\midrule
FL presence & 0.0933$^{***}$ & 0.0382$^{**}$ & 0.0760$^{***}$ & $-$0.160$^{***}$ \\
            & (0.0255)        & (0.0186)        & (0.0210)         & (---)             \\
\midrule
Item~$\times$~Year FE & YES & YES & YES & YES \\
PBU FE                & YES & YES & YES & YES \\
Observations          & 546{,}549 & 1{,}107{,}852 & --- & --- \\
\bottomrule
\end{tabular}
\begin{tablenotes}
\small
\item \textit{Notes:} Outcome is $\log$ negotiated price.
``$n_{\text{genuine}}$'' is the number of non-FL bidders in a tender;
the convite modality (Lei~8.666 Art.~22) requires at least three
bidders for the call to clear. Columns~(1)--(2) reproduce the
modality-stratified specifications of Table~\ref{tab:prices}.
Columns~(3)--(4) decompose the convite estimate by whether the
minimum-bidder constraint binds: when $n_{\text{genuine}} < 3$
(column~4), regulatory cover bidding is statutorily required; when
$n_{\text{genuine}} \geq 3$ (column~3), it is voluntary. The
voluntary subset isolates the strategic component of FL deployment.
The 7.6\% premium under voluntary cover bidding, against a $-16.0$\%
coefficient under binding constraint, is consistent with cartel-driven
deployment in the voluntary cell and with mechanical
procedural-compliance deployment in the binding cell, where FL
participation is the statutory requirement rather than a strategic
choice. Standard errors clustered at the item level.
$^{*}$~$p<0.10$, $^{**}$~$p<0.05$, $^{***}$~$p<0.01$.
\end{tablenotes}
\end{threeparttable}
\end{table}
